\documentclass[11pt,a4paper]{article}

% --------------------------------------------------------------------
%  overlap-calculator User Manual --- LaTeX source
%  Mirrors docs/MANUAL.md.  Compile with:
%      pdflatex MANUAL.tex
%      pdflatex MANUAL.tex      % second pass for TOC / refs
% --------------------------------------------------------------------

\usepackage{fontspec}
\setmainfont{Latin Modern Roman}
\setsansfont{Latin Modern Sans}
\setmonofont{Latin Modern Mono}
\usepackage[a4paper,margin=2.5cm]{geometry}
\usepackage{microtype}
\usepackage{amsmath}
\usepackage{amssymb}
\usepackage{booktabs}
\usepackage{longtable}
\usepackage{array}
\usepackage{enumitem}
\usepackage{xcolor}
\usepackage{listings}
\usepackage{seqsplit}
\usepackage{ragged2e}
\usepackage{fancyhdr}
\setlength{\headheight}{26pt}
\addtolength{\topmargin}{-12pt}
\usepackage{titlesec}
\usepackage[bookmarks=true,bookmarksopen=true,colorlinks=true,
            linkcolor=Teal,citecolor=Teal,urlcolor=Teal]{hyperref}

% -----------------------------  colours  ----------------------------
\definecolor{Teal}{RGB}{0,102,102}
\definecolor{CodeBg}{RGB}{245,245,245}
\definecolor{CodeKeyword}{RGB}{0,60,120}
\definecolor{CodeString}{RGB}{120,40,20}
\definecolor{CodeComment}{RGB}{100,100,100}
\definecolor{RuleGrey}{RGB}{180,180,180}

% Helper macro: typewriter identifier with soft-break points between every
% character, so long names like GAUSS_REFERENCE_CONCENTRATION_MOLAR can wrap
% inside narrow longtable columns without overflowing.
\newcommand{\idtt}[1]{\texttt{\seqsplit{#1}}}

% -----------------------------  listings  ---------------------------
\lstdefinelanguage{jsonish}{
  morestring=[b]",
  morecomment=[l]{//},
  sensitive=true,
}
\lstset{
  basicstyle=\ttfamily\footnotesize,
  backgroundcolor=\color{CodeBg},
  frame=single,
  framerule=0pt,
  rulecolor=\color{CodeBg},
  framesep=8pt,
  framexleftmargin=8pt,
  framexrightmargin=8pt,
  xleftmargin=8pt,
  xrightmargin=8pt,
  linewidth=\linewidth,
  aboveskip=8pt,
  belowskip=8pt,
  breaklines=true,
  breakatwhitespace=false,
  breakindent=0pt,
  postbreak=\mbox{\textcolor{gray}{$\hookrightarrow$}\space},
  columns=fullflexible,
  keepspaces=true,
  keywordstyle=\color{CodeKeyword}\bfseries,
  stringstyle=\color{CodeString},
  commentstyle=\color{CodeComment}\itshape,
  showstringspaces=false,
  tabsize=2,
  upquote=true,
}

% -----------------------------  headings  ---------------------------
\titleformat{\section}{\Large\bfseries\color{Teal}}{\thesection}{0.75em}{}
\titleformat{\subsection}{\large\bfseries}{\thesubsection}{0.6em}{}
\titleformat{\subsubsection}{\normalsize\bfseries\itshape}{\thesubsubsection}{0.5em}{}

% -----------------------------  page style  -------------------------
\pagestyle{fancy}
\fancyhf{}
\fancyhead[L]{\small overlap-calculator \textemdash{} User Manual}
\fancyhead[R]{\small \thepage}
\renewcommand{\headrulewidth}{0.4pt}
\renewcommand{\headrule}{\hbox to\headwidth{\color{RuleGrey}\leaders\hrule height \headrulewidth\hfill}}

\setlist[itemize]{leftmargin=1.2em,itemsep=2pt,topsep=2pt}
\setlist[enumerate]{leftmargin=1.4em,itemsep=2pt,topsep=2pt}

% -----------------------------  metadata  ---------------------------
\title{\vspace{-1.5cm}%
  \Huge\bfseries overlap-calculator\\[4pt]
  \Large\mdseries User Manual}
\author{Pınar Seyitdanlıoğlu}
\date{Version 1.1.0 \quad\textbullet\quad 2026-06-05}

% ====================================================================
\begin{document}
\maketitle
\thispagestyle{empty}

\vspace{1em}
\noindent\textbf{overlap-calculator} is a Python package for
reproducible batch analysis of the spectral overlap between a molecular
absorption spectrum and a reference light source. It treats two input
families as equal first-class citizens in the same run:

\begin{itemize}
  \item \textbf{Theoretical} --- Gaussian TD-DFT
        \texttt{.out}\,/\,\texttt{.log} files.
  \item \textbf{Experimental} --- tabular spectra as \textbf{CSV}
        (\texttt{.csv}), \textbf{Excel workbook} (\texttt{.xlsx}), or
        \textbf{legacy Excel} (\texttt{.xls}), with multi-sheet and
        multi-series selection.
\end{itemize}

\noindent Every TD-DFT spectrum is broadened with both Gaussian and
Lorentzian line shapes, and the tool reports per-light-source
absorbed-flux, absorbed-fraction, and shape-overlap metrics together
with publication-grade TIFF plots.

\vspace{0.5em}
\noindent This manual covers installation, the CLI, the HTTP API, the
input and output schemas, the physical assumptions and exact formulas,
and environment configuration. For worked examples on real data, see
the \texttt{case\_studies/} directory of the repository.

\vspace{1.5em}
\hrule
\vspace{0.5em}
\begin{flushleft}
{\small Project home: \url{https://github.com/pinarsyt/overlap-calculator} \\
Archived release (all versions): \url{https://doi.org/10.5281/zenodo.19944515} \\
Licensed under the GNU General Public License v3.0 or later.}
\end{flushleft}
\vspace{0.5em}
\hrule

\clearpage
\tableofcontents
\clearpage

% ====================================================================
\section{Overview and Workflow}

\begin{lstlisting}
 Gaussian TD-DFT (.out)  -->  excited states (E, lambda, f)
                                    |
                                    v                    +--- Gaussian broadening ---+
 Experimental (.csv/.xlsx) -->  A(lambda) --> eps, alpha  |                            |
                                    |                    +--- Lorentzian broadening --+
                                    v
                        x each reference light source
                                    |
                                    v
           absorbed_flux, absorbed_fraction, shape_overlap
                                    |
                                    v
         results.{csv,json,xlsx}, descriptor_summary.*, skipped_inputs.*,
         absorption / light_source / overlap / overlays TIFF plots
\end{lstlisting}

\noindent Pipeline stages:
\begin{enumerate}
  \item \textbf{Input resolution} --- A JSON manifest can be supplied
        directly or auto-generated from a directory of raw files.
        Experimental inputs expand to one row per numeric series column.
  \item \textbf{Parsing} --- Regex-based TD-DFT parser for \texttt{.out}
        files; pandas-based wavelength / signal loader for tabular
        spectra.
  \item \textbf{Spectrum reconstruction} --- Molar extinction
        \(\varepsilon(\lambda)\) from TD-DFT states via
        oscillator-strength weighting on a shared wavelength grid;
        experimental signals are treated as absorbance directly.
  \item \textbf{Beer--Lambert} ---
        \(\alpha(\lambda) = 1 - 10^{-A(\lambda)}\) with configurable
        reference concentration and path length.
  \item \textbf{Overlap metrics} --- Per
        \((\text{sample, light source, broadening})\) triplet.
  \item \textbf{Export} --- Tabular (CSV / JSON / XLSX), descriptor
        summary, skip report, and TIFF plots.
\end{enumerate}

% ====================================================================
\section{Installation}

\subsection*{Requirements}
\begin{itemize}
  \item Python \textbf{3.12.x}
  \item Optional: Conda (recommended) or Docker
\end{itemize}

\subsection*{Conda (recommended)}
\begin{lstlisting}[language=bash]
conda env create -f environment.yml
conda activate overlap-calculator
pip install -e .
\end{lstlisting}

\subsection*{venv / pip}
\begin{lstlisting}[language=bash]
python -m venv .venv
source .venv/bin/activate          # Linux / macOS
# Windows PowerShell: .venv\Scripts\Activate.ps1
# Windows CMD:        .venv\Scripts\activate.bat
pip install -e .
\end{lstlisting}

\subsection*{Docker}
The Docker image exposes the Flask API on port \texttt{8000}
\emph{inside} the container. To keep it from clashing with a native
Flask instance on the host (which uses port \texttt{8000} directly),
publish it on host port \texttt{8080}:
\begin{lstlisting}[language=bash]
docker build -t overlap-calculator .
docker run --rm -p 8080:8000 overlap-calculator
curl http://localhost:8080/health
\end{lstlisting}

\subsection*{Docker Compose}
For a managed runtime (restart policy, healthcheck, pre-populated
\texttt{GAUSS\_*} environment), use the bundled compose file. It
uses the same \textbf{host port \texttt{8080}} mapping as
\texttt{docker run} above, so the native instance on \texttt{8000}
and either docker variant on \texttt{8080} can run side by side:

\begin{lstlisting}[language=bash]
docker compose up -d            # starts the `overlap-calculator` service
curl http://localhost:8080/health
docker compose down
\end{lstlisting}

\noindent The compose service uses a Python-stdlib healthcheck
against \texttt{/health} and restarts unless explicitly stopped. All
runtime settings in \S\ref{sec:config} can be overridden via the
\texttt{environment:} block or an adjacent \texttt{.env} file.

\subsection*{Verify the install}
\begin{lstlisting}[language=bash]
overlap-calculator version
overlap-calculator --help
\end{lstlisting}

% ====================================================================
\section{Quick Start}

Drop any combination of the supported input files into
\texttt{input/files/} --- TD-DFT outputs, CSV spectra, and Excel
workbooks can live side by side in the same directory:

\begin{lstlisting}
input/files/
|-- slurm-5473089.out        # Gaussian TD-DFT output (theoretical)
|-- dye_panel.xlsx           # Excel workbook, one or more sheets (experimental)
`-- uv_vis_run3.csv          # CSV spectrum (experimental)
\end{lstlisting}

\noindent Then run the two-step pipeline:

\begin{lstlisting}[language=bash]
overlap-calculator generate-input --files-dir input/files --out input/input.json
overlap-calculator analyze --input input/input.json --out output
\end{lstlisting}

Results land under \texttt{output/tables/} and \texttt{output/plots/}
(see \S\ref{sec:outputs}).

% ====================================================================
\section{Command-Line Interface}
\label{sec:cli}

The package installs a single entry point, \texttt{overlap-calculator},
with three subcommands.

\subsection{\texttt{generate-input}}

\begin{lstlisting}[language=bash]
overlap-calculator generate-input --files-dir PATH [--out PATH] [--log-level LEVEL] [--log-format FORMAT]
\end{lstlisting}

\begin{longtable}{@{}p{3.4cm}p{3.6cm}p{7.7cm}@{}}
\toprule
\textbf{Flag} & \textbf{Default} & \textbf{Description} \\
\midrule
\endhead
\texttt{-{}-files-dir}  & \texttt{input/files}      & Directory containing TD-DFT and experimental files \\
\texttt{-{}-out}        & \texttt{input/input.json} & Output manifest path \\
\texttt{-{}-log-level}  & \texttt{INFO}             & \texttt{DEBUG}, \texttt{INFO}, \texttt{WARNING}, \texttt{ERROR}, \texttt{CRITICAL} \\
\texttt{-{}-log-format} & \texttt{text}             & \texttt{text} or \texttt{json} \\
\bottomrule
\end{longtable}

\subsection{\texttt{analyze}}

\begin{lstlisting}[language=bash]
overlap-calculator analyze --input PATH [--out PATH] [--sigma-ev FLOAT] [--wl-min FLOAT] [--wl-max FLOAT] [--num-points INT] [--concentration-m FLOAT] [--path-cm FLOAT] [--prefactor-mode {constant,frequency-resolved}] [--sigma-mode {fixed,marcus-hush}] [--reorganization-ev FLOAT] [--temperature-k FLOAT] [--calibration PATH] [--default-light-sources STR] [--light-source-file PATH]... [--plot-outputs / --no-plot-outputs] [--plot-dpi INT] [--ranking-outputs / --no-ranking-outputs] [--log-level LEVEL] [--log-format FORMAT]
\end{lstlisting}

\begin{longtable}{@{}>{\RaggedRight}p{5.0cm}>{\RaggedRight}p{3.6cm}>{\RaggedRight}p{6.1cm}@{}}
\toprule
\textbf{Flag} & \textbf{Default} & \textbf{Description} \\
\midrule
\endhead
\idtt{-{}-input}                 & \texttt{input/input.json}         & Manifest JSON (see \S\ref{sec:input}) \\
\idtt{-{}-out}                   & \texttt{output}                   & Output directory (created if absent) \\
\idtt{-{}-sigma-ev}              & \(0.30\)                          & Broadening standard deviation, eV \\
\idtt{-{}-wl-min}                & \(200.0\)                         & Lower bound of the wavelength grid, nm \\
\idtt{-{}-wl-max}                & \(800.0\)                         & Upper bound of the wavelength grid, nm \\
\idtt{-{}-num-points}            & \(10000\)                         & Number of grid points \\
\idtt{-{}-concentration-m}       & \(1\text{e-}5\)                   & Beer--Lambert reference concentration, mol\,L\(^{-1}\) \\
\idtt{-{}-path-cm}               & \(1.0\)                           & Beer--Lambert reference path length, cm \\
\idtt{-{}-prefactor-mode}        & \texttt{constant}                 & \texttt{constant} (integrated intensity; default) or \texttt{frequency-resolved} (keeps the wavenumber factor inside the integrand; relevant for NIR bands). See \S\ref{sec:prefactor}. \\
\idtt{-{}-sigma-mode}            & \texttt{fixed}                    & \texttt{fixed} (use \idtt{-{}-sigma-ev}) or \texttt{marcus-hush} (\(\sigma=\sqrt{2\lambda k_B T}\) from \idtt{-{}-reorganization-ev} and \idtt{-{}-temperature-k}). See \S\ref{sec:prefactor}. \\
\idtt{-{}-reorganization-ev}     & \(0.30\)                          & Reorganization energy \(\lambda\), eV; used when \texttt{-{}-sigma-mode marcus-hush} \\
\idtt{-{}-temperature-k}         & \(298.15\)                        & Temperature in K for the Marcus--Hush width \\
\idtt{-{}-calibration}           & \emph{none}                       & Optional TOML calibration file (CLI and library only). See \S\ref{sec:calibration}. \\
\idtt{-{}-default-light-sources} & \texttt{AM15G,LEDB4,\ldots}       & Comma-separated bundled-source list \\
\idtt{-{}-light-source-file}     & \emph{none}                       & Repeatable path to a custom light-source CSV \\
\idtt{-{}-plot-outputs}          & on                                & Emit TIFF plots (\idtt{-{}-no-plot-outputs} to disable) \\
\idtt{-{}-plot-dpi}              & \(400\)                           & TIFF plot resolution; use \(600\) for larger publication exports \\
\idtt{-{}-ranking-outputs}       & on                                & Emit grouped ranking tables and bar charts (gaussian/lorentzian \(\times\) \texttt{absorbed\_fraction} / \texttt{shape\_overlap}); disable with \idtt{-{}-no-ranking-outputs} \\
\idtt{-{}-log-level}             & \texttt{INFO}                     & Logging verbosity \\
\idtt{-{}-log-format}            & \texttt{text}                     & \texttt{text} or \texttt{json} \\
\bottomrule
\end{longtable}

\subsection{\texttt{version}}
Prints the installed package version.

\begin{lstlisting}[language=bash]
overlap-calculator version
\end{lstlisting}

% ====================================================================
\section{HTTP API}
\label{sec:api}

\noindent The examples in this section use port \texttt{8000}, which
is the \textbf{native Flask} port (started with \texttt{python -m
overlap\_calculator.api.app}). If you run the same API under
\textbf{Docker} (either \texttt{docker run -p 8080:8000} or
\texttt{docker compose up}), use port \texttt{8080} instead --- the
request bodies are otherwise identical.

\subsection{Endpoints}

\begin{longtable}{@{}p{1.8cm}p{2.8cm}p{10.6cm}@{}}
\toprule
\textbf{Method} & \textbf{Path} & \textbf{Purpose} \\
\midrule
\endhead
GET  & \texttt{/health}  & Liveness probe \textemdash{} returns \texttt{\{"status":"ok"\}} \\
POST & \texttt{/analyze} & Multipart upload; returns a ZIP of results \\
\bottomrule
\end{longtable}

\noindent \texttt{POST /analyze} accepts the following form fields:

\begin{longtable}{@{}>{\RaggedRight}p{5.0cm}>{\RaggedRight}p{3.0cm}>{\RaggedRight}p{6.7cm}@{}}
\toprule
\textbf{Form field} & \textbf{Type} & \textbf{Description} \\
\midrule
\endhead
\idtt{files}                  & \texttt{file[]} (req.) & Mixed \texttt{.out}/\texttt{.csv}/\texttt{.xlsx}/\texttt{.xls} inputs \\
\idtt{light\_source\_files}   & \texttt{file[]} (opt.) & Custom light-source CSVs \\
\idtt{light\_source\_csv}     & \texttt{file} (opt.)   & Back-compat alias for a single custom CSV \\
\idtt{sigma\_ev}              & float (opt.)           & Broadening \(\sigma\), eV (default \(0.30\)) \\
\idtt{wl\_min}, \idtt{wl\_max} & float (opt.)        & Wavelength window bounds, nm \\
\idtt{num\_points}            & int (opt.)             & Grid size (default \(10000\)) \\
\idtt{concentration\_m}       & float (opt.)           & Beer--Lambert reference concentration (default \(1\text{e-}5\)) \\
\idtt{path\_cm}               & float (opt.)           & Beer--Lambert reference path length (default \(1.0\)) \\
\idtt{prefactor\_mode}        & str (opt.)             & \texttt{constant} (default) or \texttt{frequency-resolved} \\
\idtt{sigma\_mode}            & str (opt.)             & \texttt{fixed} (default) or \texttt{marcus-hush} \\
\idtt{reorganization\_ev}     & float (opt.)           & Reorganization energy \(\lambda\), eV; used when \texttt{sigma\_mode=marcus-hush} (default \(0.30\)) \\
\idtt{temperature\_k}         & float (opt.)           & Temperature in K for the Marcus--Hush width (default \(298.15\)) \\
\idtt{default\_light\_sources}& str (opt.)             & Comma-separated bundled-source list \\
\idtt{plot\_outputs}          & bool (opt.)            & Enable / disable plot generation \\
\idtt{plot\_dpi}              & int (opt.)             & TIFF plot resolution (default \(400\); use \(600\) for larger exports) \\
\idtt{ranking\_outputs}       & bool (opt.)            & Enable / disable grouped ranking tables and bar charts (default \texttt{true}) \\
\idtt{sheet\_overrides}       & str (opt.)             & JSON object \idtt{\{"<original-filename>":"<sheet-name>"\}} forcing a specific sheet for uploaded \texttt{.xlsx}/\texttt{.xls} workbooks \\
\bottomrule
\end{longtable}

\noindent Response: \texttt{application/zip} containing every file
\texttt{overlap-calculator analyze} would have written. Uploads are
processed in a per-request workspace and are not persisted on the
server.

\subsection{Example}
\begin{lstlisting}[language=bash]
curl -X POST http://localhost:8000/analyze \
  -F "files=@input/files/theoretical.out" \
  -F "files=@input/files/experimental.xlsx" \
  -F "sigma_ev=0.25" \
  -F "plot_dpi=400" \
  -F "ranking_outputs=true" \
  -F "default_light_sources=AM15G,LEDB4" \
  --output analysis_outputs.zip
\end{lstlisting}

\noindent To force a specific sheet in a multi-sheet workbook, add
a \idtt{sheet\_overrides} form field whose value is a JSON object keyed
by the original filename:

\begin{lstlisting}[language=bash]
curl -X POST http://localhost:8000/analyze \
  -F "files=@input/files/dye_panel.xlsx" \
  -F 'sheet_overrides={"dye_panel.xlsx":"raw"}' \
  --output analysis_outputs.zip
\end{lstlisting}

\noindent Filenames that do not appear in the map use the same
auto-discovery as \texttt{generate-input}: the first sheet with a
wavelength column and at least one numeric signal column is selected.

A Postman collection is provided at
\texttt{docs/postman\_collection.json}.

% ====================================================================
\section{Input Format}
\label{sec:input}

\subsection{Accepted input files}

\texttt{overlap-calculator} accepts two input families in the same run.
Both are first-class: a single manifest can mix them freely, and the
pipeline produces identical downstream columns, plots, and overlap
metrics.

\begin{longtable}{@{}>{\RaggedRight}p{2.6cm}>{\RaggedRight}p{2.8cm}>{\RaggedRight}p{4.2cm}>{\RaggedRight}p{4.7cm}@{}}
\toprule
\textbf{Family} & \textbf{Extensions} & \textbf{Parser} & \textbf{Notes} \\
\midrule
\endhead
Theoretical  & \texttt{.out}, \texttt{.log}           & Regex-based Gaussian TD-DFT parser                     & File must contain a TD route and \texttt{Excited State} lines \\
Experimental & \texttt{.csv}                          & \texttt{pandas.read\_csv}                              & One wavelength column and one or more numeric signal columns \\
Experimental & \texttt{.xlsx}, \texttt{.xls}          & \texttt{pandas.read\_excel} (\texttt{openpyxl}/\texttt{xlrd}) & Multi-sheet workbooks supported via \texttt{sheet\_name}; \texttt{generate-input} auto-selects the first analyzable sheet \\
\bottomrule
\end{longtable}

\noindent Bulk uploads via the HTTP API share the same matrix (see
\idtt{GAUSS\_UPLOAD\_ALLOWED\_EXTS} in \S\ref{sec:config}).

\subsection{Manifest (\texttt{input.json})}

\begin{lstlisting}[language=jsonish]
[
  {
    "input_type": "theoretical",
    "sample_id": "B1",
    "source_path": "files/slurm-5473089.out",
    "series_name": null,
    "sheet_name": null
  },
  {
    "input_type": "experimental",
    "sample_id": "channel_A",
    "source_path": "files/dye_panel.xlsx",
    "series_name": "channel_A",
    "sheet_name": "raw"
  }
]
\end{lstlisting}

\begin{longtable}{@{}p{2.8cm}p{1.6cm}p{10.15cm}@{}}
\toprule
\textbf{Field} & \textbf{Req?} & \textbf{Description} \\
\midrule
\endhead
\texttt{input\_type}   & yes   & \texttt{"theoretical"} (TD-DFT \texttt{.out}) or \texttt{"experimental"} (tabular spectrum) \\
\texttt{sample\_id}    & yes   & Unique sample identifier; used in filenames and result rows \\
\texttt{source\_path}  & yes   & Path to the raw file; resolved relative to the manifest, then to the CWD \\
\texttt{series\_name}  & exp.  & Column name of the numeric series in a tabular file \\
\texttt{sheet\_name}   & exp.  & Sheet name for \texttt{.xlsx}/\texttt{.xls} inputs; in hand-written manifests, omit or \texttt{null} to read the first workbook sheet \\
\bottomrule
\end{longtable}

\noindent \texttt{generate-input} fills these automatically. The
\texttt{sample\_id} is derived as follows:

\begin{itemize}
  \item \textbf{Theoretical} (\texttt{.out}) --- if the Gaussian route
        contains a \texttt{\%chk=<name>.chk} line, the checkpoint stem
        is used (e.g.\ \texttt{\%chk=B1\_td.chk} \(\Rightarrow\)
        \texttt{sample\_id = "B1\_td"}). This keeps human-meaningful
        molecule labels even when the raw files are queue-system
        artefacts such as \texttt{slurm-5473089.out}. If no
        \texttt{\%chk} line is present, the file stem is used as a
        fallback.
  \item \textbf{Experimental} --- one manifest entry is produced per
        numeric series column, with
        \texttt{sample\_id = "<series\_name>"} (the series column name
        itself, so \texttt{Coumarin\_1} rather than a long file-stem
        prefix).
\end{itemize}

\noindent If a manifest is supplied by hand and \texttt{sample\_id} is
omitted, the same resolution rule is applied at analysis time.

\subsection{Experimental data template}

The same column layout applies to CSV, XLSX, and XLS files.

\paragraph{CSV (\texttt{.csv}).}
\begin{lstlisting}
wavelength_nm,sample_A,sample_B
300,0.12,0.08
310,0.15,0.10
320,0.18,0.14
\end{lstlisting}

\paragraph{Excel workbook (\texttt{.xlsx}\,/\,\texttt{.xls}).}
The first row is interpreted as the header. Any numeric columns beside
the wavelength axis are treated as independent samples:

\begin{longtable}{@{}p{3.2cm}p{2.6cm}p{2.6cm}p{2.6cm}@{}}
\toprule
\textbf{wavelength\_nm} & \textbf{channel\_A} & \textbf{channel\_B} & \textbf{channel\_C} \\
\midrule
\endhead
300 & 0.12 & 0.08 & 0.05 \\
310 & 0.15 & 0.10 & 0.07 \\
320 & 0.18 & 0.14 & 0.09 \\
\bottomrule
\end{longtable}

\noindent Multi-sheet workbooks are fully supported. During
\texttt{generate-input}, if no sheet override is supplied, the workbook
sheets are scanned in order and the first sheet with a wavelength column
and at least one numeric signal column is selected. In a hand-written
manifest, if \texttt{sheet\_name} is omitted from an entry (or set to
\texttt{null}), analysis reads the first workbook sheet. To force a
specific sheet, edit the generated manifest entry or author it by hand:

\begin{lstlisting}[language=jsonish]
{
  "input_type": "experimental",
  "sample_id": "dye_panel__channel_A",
  "source_path": "files/dye_panel.xlsx",
  "series_name": "channel_A",
  "sheet_name": "raw"
}
\end{lstlisting}

\paragraph{Shared requirements (both formats).}
\begin{itemize}
  \item Exactly one wavelength column; accepted aliases
        (case-insensitive): \texttt{wavelength},
        \idtt{wavelength\_nm}, \texttt{lambda}, \texttt{wl},
        \texttt{nm}. If no named wavelength column is found, the first
        numeric column is used.
  \item One or more numeric signal columns. Each signal is interpreted
        as absorbance \(A(\lambda)\) and fed directly to the
        Beer--Lambert absorptance definition
        (\S\ref{sec:methodology}).
  \item Rows with non-numeric or missing values in either the
        wavelength or signal column are dropped silently; a series must
        retain at least 10 valid rows to be analysed.
\end{itemize}

\subsection{Custom light sources}
A custom light source is a two-column CSV with one wavelength column
and a single numeric intensity column. The filename stem becomes the
\texttt{light\_source\_name} in the results table.

% ====================================================================
\section{Output Schema}
\label{sec:outputs}

\begin{lstlisting}
output/
|-- run_manifest.json                                          % provenance record
|-- tables/
|   |-- results.{csv,json,xlsx}
|   |-- results_timings.{csv,json,xlsx}
|   |-- descriptor_summary.{csv,xlsx}
|   |-- ranking_by_light_source__<metric>.{csv,json,xlsx}    % 4 metric variants
|   |-- ranking_by_sample__<metric>.{csv,json,xlsx}           % 4 metric variants
|   `-- skipped_inputs.{csv,json}           % emitted only if any inputs were skipped
`-- plots/
    |-- absorption/<sample>__{absolute,normalized}.tiff
    |-- light_source/<light>.tiff
    |-- overlap/<sample>__<light>__combined__{absolute,normalized}.tiff
    |-- overlap/<sample>__<light>__<method>__{absolute,normalized}.tiff
    |-- overlays/absorptance_overlay__<method>__{absolute,normalized}.tiff
    |-- ranking/by_light_source/<metric>/<light>.tiff         % bar chart per light source
    `-- ranking/by_sample/<metric>/<sample>.tiff              % bar chart per sample
\end{lstlisting}

\noindent \texttt{<metric>} cycles through the four overlap descriptors:
\idtt{gaussian\_absorbed\_fraction},
\idtt{lorentzian\_absorbed\_fraction},
\idtt{gaussian\_shape\_overlap},
\idtt{lorentzian\_shape\_overlap}. The ranking tables and bar
charts can be suppressed with \idtt{-{}-no-ranking-outputs} (CLI)
or \idtt{ranking\_outputs=false} (API).

\subsection{\texttt{results.*} columns}
One row per \((\text{sample\_id},\text{light\_source\_name})\) pair.
Shared columns:

\begin{longtable}{@{}>{\RaggedRight}p{6.2cm}>{\RaggedRight}p{9.2cm}@{}}
\toprule
\textbf{Column} & \textbf{Description} \\
\midrule
\endhead
\idtt{source\_type}                    & \texttt{theoretical} or \texttt{experimental} \\
\idtt{sample\_id}                      & Unique sample identifier \\
\idtt{light\_source\_name}             & \texttt{AM15G}, \texttt{LEDB1}, \texttt{LEDB2}, \texttt{LEDB3}, \texttt{LEDB4}, \texttt{CIEFL10}, or custom stem \\
\idtt{sigma\_ev}                       & Broadening \(\sigma\) used, eV (theoretical inputs) \\
\idtt{light\_flux\_total}              & \(\int I(\lambda)\,d\lambda\) over the working grid \\
\idtt{reference\_concentration\_molar} & Beer--Lambert \(c\) used, mol\,L\(^{-1}\) \\
\idtt{reference\_path\_cm}             & Beer--Lambert \(L\) used, cm \\
\idtt{absorption\_unit}                & Unit of \(\varepsilon\) / \(A\) in the exported row \\
\idtt{light\_source\_unit}             & Unit of \(I(\lambda)\) (e.g.\ \texttt{W\,m\(^{-2}\)\,nm\(^{-1}\)} for AM1.5G) \\
\idtt{absorbed\_flux\_unit}            & Unit of the absorbed-flux integral \\
\idtt{td\_method\_basis}               & Route metadata from the TD-DFT file (null for experimental) \\
\idtt{td\_solvent}                     & PCM solvent from the TD-DFT file (null for experimental) \\
\idtt{prefactor\_mode}                 & \texttt{constant} or \texttt{frequency-resolved} --- the prefactor convention used for this row \\
\idtt{sigma\_mode}                     & \texttt{fixed} or \texttt{marcus-hush} --- the width-source used for this row \\
\idtt{software\_version}               & Package version string at run time \\
\bottomrule
\end{longtable}

\noindent Per-broadening columns (prefixes \texttt{gaussian\_} and
\texttt{lorentzian\_}):

\begin{longtable}{@{}>{\RaggedRight}p{4.8cm}>{\RaggedRight}p{10.6cm}@{}}
\toprule
\textbf{Suffix} & \textbf{Meaning} \\
\midrule
\endhead
\idtt{lambda\_max\_nm}               & Peak wavelength of the broadened absorptance \\
\idtt{molar\_absorptivity\_max}      & Peak \(\varepsilon\) in M\(^{-1}\,\)cm\(^{-1}\) (theoretical only) \\
\idtt{absorptance\_max}              & Peak Beer--Lambert \(\alpha(\lambda)\) \\
\idtt{absorbed\_flux}                & \(\int \alpha(\lambda)\,I(\lambda)\,d\lambda\) \\
\idtt{absorbed\_fraction}            & \(\text{absorbed\_flux} / \text{light\_flux\_total}\in[0,1]\) \\
\idtt{shape\_overlap}                & Dimensionless shape comparator on max-normalised spectra \\
\bottomrule
\end{longtable}

\noindent Broadening deltas (Gaussian \(-\) Lorentzian):
\idtt{delta\_absorbed\_flux},
\idtt{delta\_absorbed\_fraction},
\idtt{delta\_shape\_overlap},
\idtt{abs\_delta\_absorbed\_fraction}.

\subsection{\texttt{descriptor\_summary.*}}
A condensed table of the most actionable per-sample descriptors across
all light sources, suitable as a starting point for ranking.

\subsection{\texttt{ranking\_by\_light\_source\_\_<metric>.*} and \texttt{ranking\_by\_sample\_\_<metric>.*}}
\label{sec:rankings}

Two complementary grouped rankings answer the two practical questions
that arise from a multi-sample, multi-light-source run:

\begin{itemize}
\item \textbf{\idtt{ranking\_by\_light\_source\_\_<metric>.\{csv,json,xlsx\}}}
      \textemdash{} within each light source, samples are ordered from
      the highest to the lowest value of the chosen overlap metric. Use
      this to ask: \emph{``Given this light source, which molecule
      absorbs (or matches in shape) best?''}
\item \textbf{\idtt{ranking\_by\_sample\_\_<metric>.\{csv,json,xlsx\}}}
      \textemdash{} within each sample, light sources are ordered the
      same way. Use this to ask: \emph{``For this molecule, which light
      source gives the strongest overlap?''}
\end{itemize}

\noindent Each table is emitted four times \textemdash{} once per
\texttt{<metric>} \textemdash{} covering both broadening methods
(\texttt{gaussian\_*}, \texttt{lorentzian\_*}) and both metric families
(\texttt{absorbed\_fraction}, \texttt{shape\_overlap}):

\begin{longtable}{@{}>{\RaggedRight}p{6.4cm}>{\RaggedRight}p{9.0cm}@{}}
\toprule
\textbf{\texttt{<metric>}} & \textbf{Sorts by} \\
\midrule
\endhead
\idtt{gaussian\_absorbed\_fraction}    & Beer--Lambert absorbed fraction with Gaussian broadening; intensity-weighted \\
\idtt{lorentzian\_absorbed\_fraction}  & Beer--Lambert absorbed fraction with Lorentzian broadening; intensity-weighted \\
\idtt{gaussian\_shape\_overlap}        & Shape-only correlation (intensity-independent) under Gaussian broadening \\
\idtt{lorentzian\_shape\_overlap}      & Shape-only correlation (intensity-independent) under Lorentzian broadening \\
\bottomrule
\end{longtable}

\noindent Both interpretations are reported because they answer
different scientific questions: \idtt{absorbed\_fraction} tracks the
fraction of incident photons that are actually absorbed under the
chosen Beer--Lambert reference \((c, L)\) and is the right figure of
merit for photocatalysis, photovoltaic action, or photo-excitation
yield; \idtt{shape\_overlap} is a brightness-independent spectral
matching score, appropriate when comparing differently-scaled
illuminants on equal footing (\S\ref{sec:metric-choice}).

Each row carries a \texttt{rank} column (1 = best within the group)
and all four metric columns side by side, so the file can be
re-sorted client-side without rerunning the pipeline. Companion bar
charts are written to
\idtt{plots/ranking/by\_light\_source/<metric>/<light>.tiff} and
\idtt{plots/ranking/by\_sample/<metric>/<sample>.tiff}; titles and
axis labels embed the chosen metric and the Beer--Lambert reference
conditions to keep the figures unambiguous.

\subsection{\texttt{skipped\_inputs.*}}
Emitted whenever at least one input is rejected. The
\texttt{tables/skipped\_inputs.\{csv,json\}} rows carry
\texttt{sample\_id} and \texttt{reason} columns. The manifest-side
\texttt{input/skipped\_inputs.json} produced by \texttt{generate-input}
additionally carries the resolved \texttt{source\_path}. Each
\texttt{reason} string is prefixed with the stage at which the input
was skipped (\texttt{INPUT\_SKIPPED} / \texttt{INPUT\_ERROR} /
\texttt{ANALYSIS\_ERROR}) followed by the error message.

\subsection{\texttt{results\_timings.*}}
One row per \idtt{(sample\_id, light\_source\_name)} pair. Shared
columns: \idtt{source\_type}, \idtt{sample\_id},
\idtt{light\_source\_name}, \idtt{plot\_ms}. Each broadening method
contributes a five-column block (\idtt{gaussian\_parse\_tddft\_ms},
\idtt{gaussian\_broadening\_ms}, \idtt{gaussian\_light\_source\_ms},
\idtt{gaussian\_integrals\_ms}, \idtt{gaussian\_total\_sample\_ms},
and the matching \idtt{lorentzian\_*} quintuplet). The row ends with
\idtt{total\_pair\_ms} (sum of both methods' \idtt{total\_sample\_ms})
and \idtt{total\_with\_plot\_ms} (\idtt{total\_pair\_ms + plot\_ms}).

\subsection{Plots}
TIFF plots are emitted at 400\,dpi by default and include both absolute
and max-normalised variants. Use \idtt{-{}-plot-dpi 600} on the CLI or
\idtt{plot\_dpi=600} in the API when larger publication-resolution
files are worth the additional disk and transfer size.

\begin{longtable}{@{}p{3.8cm}p{11.4cm}@{}}
\toprule
\textbf{Directory} & \textbf{Content} \\
\midrule
\endhead
\texttt{absorption/}                                  & Per-sample absorptance spectra, both broadenings overlaid \\
\texttt{light\_source/}                               & Each reference light source in isolation \\
\texttt{overlap/<combined>}                           & (sample, light) overlap \textemdash{} both broadenings on one axes \\
\texttt{overlap/<method>}                             & (sample, light) overlap \textemdash{} single broadening \\
\texttt{overlays/}                                    & Multi-sample absorptance overlays per broadening \\
\idtt{ranking/by\_light\_source/<metric>/}            & Bar chart per light source ranking samples by \texttt{<metric>} (descending) \\
\idtt{ranking/by\_sample/<metric>/}                   & Bar chart per sample ranking light sources by \texttt{<metric>} (descending) \\
\bottomrule
\end{longtable}

\noindent The overlap plots shade the geometric overlap area
\(\min(\alpha(\lambda),\,\hat I(\lambda))\) (absolute panel) and
\(\min(\hat\alpha(\lambda),\,\hat I(\lambda))\) (normalised panel).
In the normalised panel the shaded region is exactly the integrand of
\texttt{shape\_overlap}, so the figure and the numerical metric
report the same quantity: the area bounded above by the lower of the
two curves at every wavelength.

\subsection{\texttt{run\_manifest.json}}
\label{sec:manifest}

Every \texttt{analyze} run writes \texttt{run\_manifest.json} directly
inside the output root directory (at the same level as \texttt{tables/}
and \texttt{plots/}). The file records enough information to trace any
result back to its exact inputs and parameters.

\noindent Top-level fields:

\begin{longtable}{@{}>{\RaggedRight}p{4.4cm}>{\RaggedRight}p{11.0cm}@{}}
\toprule
\textbf{Field} & \textbf{Description} \\
\midrule
\endhead
\idtt{software\_version}   & Installed package version string \\
\idtt{git\_commit}         & Full git SHA of the running commit; \texttt{null} if unavailable \\
\idtt{generated\_at\_utc}  & ISO-8601 UTC timestamp of the run \\
\idtt{parameters}          & Full resolved runtime parameter set: \texttt{sigma\_ev}, \texttt{prefactor\_mode}, \texttt{sigma\_mode}, \texttt{reorganization\_ev}, \texttt{temperature\_k}, wavelength grid, Beer--Lambert constants, light-source lists, and the full calibration block when \texttt{-{}-calibration} was supplied \\
\idtt{light\_sources}      & Sorted list of light-source names used in the run \\
\idtt{inputs}              & Array --- one object per source file (see below) \\
\idtt{results\_count}      & Number of \texttt{RunResult} rows produced \\
\idtt{skipped\_count}      & Number of inputs that were skipped \\
\bottomrule
\end{longtable}

\noindent Each element of \texttt{inputs}:

\begin{longtable}{@{}>{\RaggedRight}p{3.4cm}>{\RaggedRight}p{12.0cm}@{}}
\toprule
\textbf{Field} & \textbf{Description} \\
\midrule
\endhead
\idtt{sample\_id}    & Sample identifier from the manifest \\
\idtt{input\_type}   & \texttt{"theoretical"} or \texttt{"experimental"} \\
\idtt{source\_path}  & Run-directory-relative path to the source file (never an absolute machine path); content fingerprinted by \idtt{sha256} \\
\idtt{series\_name}  & Series column name (experimental inputs only; otherwise \texttt{null}) \\
\idtt{sheet\_name}   & Sheet name for Excel inputs with an explicit sheet override; otherwise \texttt{null} \\
\idtt{sha256}        & SHA-256 hex digest of the file contents at run time; \texttt{null} if the file was not readable \\
\bottomrule
\end{longtable}

% ====================================================================
\section{Bundled Light Sources}
\label{sec:lightsources}

\begin{longtable}{@{}p{2cm}p{8cm}p{4.35cm}@{}}
\toprule
\textbf{Name} & \textbf{Source} & \textbf{Unit} \\
\midrule
\endhead
\texttt{AM15G}   & NREL ASTM G-173 AM1.5G reference solar spectrum      & W\,m\(^{-2}\,\)nm\(^{-1}\) \\
\texttt{LEDB1}   & CIE 15 illuminant LED-B1 (opt-in; not in default set) & relative (shape only) \\
\texttt{LEDB2}   & CIE 15 illuminant LED-B2                             & relative (shape only) \\
\texttt{LEDB3}   & CIE 15 illuminant LED-B3                             & relative \\
\texttt{LEDB4}   & CIE 15 illuminant LED-B4                             & relative \\
\texttt{CIEFL10} & CIE 15 fluorescent illuminant FL10                   & relative \\
\bottomrule
\end{longtable}

\noindent Unit policy: because \texttt{AM15G} carries an absolute
spectral irradiance, \idtt{absorbed\_flux} for AM1.5G has units of
W\,m\(^{-2}\); for the relative CIE LED / FL sources
\idtt{absorbed\_flux\_unit} is \idtt{relative\_integral} and only
\idtt{absorbed\_fraction} and \idtt{shape\_overlap} are directly
comparable across dyes.

\paragraph{Third-party spectral data.}
The \texttt{LEDB1}, \texttt{LEDB2}, \texttt{LEDB3}, \texttt{LEDB4}, and
\texttt{CIEFL10} spectra are official \textbf{CIE} data sets
derived from \emph{CIE 015:2018 Colorimetry, 4th Edition}
(International Commission on Illumination,
\url{https://cie.co.at/}). The \texttt{AM15G} spectrum is the ASTM
G-173-03 AM1.5 global reference distributed by the U.S.\ National
Renewable Energy Laboratory (NREL). See \S\ref{sec:citation} for
full citations and BibTeX entries.

% ====================================================================
\section{Methodology}
\label{sec:methodology}

\subsection{Extinction spectrum from TD-DFT}

Each excited state \(i\) contributes a normalised line profile in
wavenumber space \(\tilde\nu = 10^{7}/\lambda[\text{nm}]\) (cm\(^{-1}\)):

\begin{align}
g_i(\tilde\nu) &= \frac{1}{\sigma\sqrt{2\pi}}\,
                   \exp\!\Bigl(-\tfrac{(\tilde\nu - \tilde\nu_i)^2}{2\sigma^2}\Bigr), \\[2pt]
l_i(\tilde\nu) &= \frac{1}{\pi}\,\frac{\gamma}{(\tilde\nu - \tilde\nu_i)^2 + \gamma^2},
\qquad \gamma = \sqrt{2\ln 2}\,\sigma \approx 1.1774\,\sigma,
\end{align}

with \(\sigma\) in cm\(^{-1}\) obtained from the CLI \texttt{-{}-sigma-ev}
via \(\sigma_{\text{cm}^{-1}} = \sigma_{\text{eV}} \cdot 8065.544\).
Here \(\sigma\) is the Gaussian standard deviation
(\(\mathrm{FWHM}_{\mathrm{G}} = 2\sqrt{2\ln 2}\,\sigma\approx 2.3548\,\sigma\));
the Lorentzian half-width-at-half-maximum \(\gamma\) is set to
\(\sqrt{2\ln 2}\,\sigma\) so that \emph{both line shapes share the
same \emph{FWHM}} under a given \texttt{-{}-sigma-ev} value. This
removes the implicit width mismatch between the two broadening models
and keeps the Gaussian--Lorentzian sensitivity comparison a pure
line-shape probe; the \texttt{delta\_*} columns therefore reflect
tail behaviour rather than width difference. Both profiles are
normalised in cm\(^{-1}\): \(\int p(\tilde\nu)\,d\tilde\nu = 1\). The
molar extinction coefficient is then

\begin{equation}
\varepsilon(\lambda) = 2.315\times 10^{8}\;
  \sum_i f_i\, \text{profile}_i\!\bigl(\tilde\nu(\lambda)\bigr)
  \quad [\text{M}^{-1}\,\text{cm}^{-1}],
\end{equation}

where the prefactor comes from the integrated-absorption /
oscillator-strength relation
\(\int \varepsilon\,d\tilde\nu = 2.315\times 10^{8}\,\sum_i f_i\).

\subsection{Beer--Lambert absorptance}

\begin{equation}
A(\lambda) = \varepsilon(\lambda)\,c\,L,\qquad
\alpha(\lambda) = 1 - 10^{-A(\lambda)}
\end{equation}

with defaults \(c = 1\times 10^{-5}\,\text{M}\), \(L = 1\,\text{cm}\)
(overridable via \texttt{-{}-concentration-m} / \texttt{-{}-path-cm}).
For experimental inputs the user-provided signal is treated as
absorbance \(A\) directly and fed to the same \(\alpha\) definition
without the Beer--Lambert conversion.

\subsection{Light-overlap metrics}

\begin{align}
\text{absorbed\_flux}      &= \int \alpha(\lambda)\,I(\lambda)\,d\lambda, \\
\text{light\_flux\_total}  &= \int I(\lambda)\,d\lambda, \\
\text{absorbed\_fraction}  &= \frac{\text{absorbed\_flux}}{\text{light\_flux\_total}} \in [0,1], \\
\text{shape\_overlap}      &= \frac{\displaystyle\int \min\!\bigl(\hat\alpha(\lambda),\,\hat I(\lambda)\bigr)\,d\lambda}
                                   {\displaystyle\int \hat I(\lambda)\,d\lambda},
\end{align}

where \(\hat\alpha\) and \(\hat I\) denote max-normalised shapes.
\texttt{shape\_overlap} is the classical geometric spectral-overlap
area: at every wavelength the integrand follows whichever of the two
normalised profiles is smaller, so the shaded region in a normalised
overlap plot equals the numerator directly. This descriptor is
independent of absolute strength and lets LED / FL sources (relative
intensities) still be compared.

\subsection{Choosing between \texttt{absorbed\_fraction} and
\texttt{shape\_overlap}}
\label{sec:metric-choice}

The two scalar overlap descriptors answer different questions and
are \emph{both} emitted in the ranking tables and bar charts
(\S\ref{sec:rankings}) so that no single metric is forced on the
user:

\begin{itemize}
  \item \textbf{\texttt{absorbed\_fraction}} is the right figure of
        merit when the application cares about the \emph{number of
        photons (or the energy) actually absorbed} under a defined
        Beer--Lambert reference state \((c, L)\) \textemdash{} for
        example, photocatalytic turnover under a given illuminant,
        photovoltaic action spectra, or photo-excitation yield. It
        folds in \emph{both} spectral overlap and the absolute
        spectral intensity of the light source: a brighter source
        over the absorbing band is rewarded.
  \item \textbf{\texttt{shape\_overlap}} is intensity-independent
        and answers the brightness-blind question \emph{``how well
        does the light source's spectral envelope match the
        molecule's absorption envelope?''} It is the appropriate
        descriptor when the spectra under comparison carry
        incompatible intensity units (e.g.\ \texttt{W\,m\(^{-2}\)\,nm\(^{-1}\)}
        for AM1.5G versus \texttt{relative} for the CIE LED
        templates) and the reader needs a fair side-by-side
        ordering.
\end{itemize}

\noindent Each descriptor is reported under both Gaussian and
Lorentzian broadening (\S\ref{sec:methodology}), so the four ranking
tables together expose the full sensitivity surface to the line-shape
choice.

\subsection{Skip policy}

An input is diverted to \texttt{skipped\_inputs.*} when
\begin{itemize}
  \item the \texttt{.out} file lacks a TD route or any
        \texttt{Excited State} line (tagged
        \idtt{INPUT\_SKIPPED: file is not a TD-DFT output});
  \item an experimental file cannot be read or has no numeric signal
        column (tagged
        \idtt{INPUT\_SKIPPED: failed to read experimental file} or
        \idtt{INPUT\_SKIPPED: experimental file has no numeric series});
  \item parsing fails at analysis time (tagged
        \idtt{ANALYSIS\_ERROR: Failed to parse \ldots}).
\end{itemize}

\subsection{Choice of prefactor and broadening width}
\label{sec:prefactor}

\subsubsection*{Prefactor modes}

The oscillator-strength prefactor
\(P = 2.315 \times 10^8\,\text{M}^{-1}\,\text{cm}^{-2}\) comes from
the Hilborn/Mulliken oscillator-strength relation: the integrated band
area equals \(P \cdot f\).

\paragraph{Constant prefactor (\texttt{-{}-prefactor-mode constant}, default).}
\begin{equation}
\varepsilon(\nu) = \sum_i P f_i\, g_i(\nu - \nu_i)
\end{equation}
where \(g_i\) is a normalized lineshape. The band area equals
\(P f_i\) independent of \(\nu_i\). This is the validated convention
for the 200--800\,nm visible range.

\paragraph{Frequency-resolved prefactor (\texttt{-{}-prefactor-mode frequency-resolved}).}
\begin{equation}
\varepsilon(\nu) = \sum_i P f_i \frac{\nu}{\nu_i}\, g_i(\nu - \nu_i)
\end{equation}
The integrated band area is unchanged for a symmetric lineshape. The two
modes differ only in peak height (by approximately
\(+\tfrac{1}{2}(\sigma/\nu_i)^2\)) and in a small blue-shift of the peak,
of order \(O((\sigma/\nu_i)^2)\). This difference is below \(\sim\)2\,\%
over 200--800\,nm for \(\sigma = 0.30\)\,eV, and grows into the NIR.

\textbf{Recommendation:} use \texttt{constant} for the validated visible
range; use \texttt{frequency-resolved} for NIR bands.

\subsubsection*{Lineshape definitions}

Gaussian:
\begin{equation}
g(x) = \frac{1}{\sigma_{\text{cm}}\sqrt{2\pi}}\,
        \exp\!\Bigl(-\frac{x^2}{2\sigma_{\text{cm}}^2}\Bigr),
\qquad \sigma_{\text{cm}} = \sigma_{\text{eV}} \times 8065.5439\;\text{cm}^{-1}.
\end{equation}

Lorentzian (FWHM-matched to the Gaussian):
\begin{equation}
l(x) = \frac{\gamma/\pi}{x^2 + \gamma^2},\qquad
\gamma = \sigma_{\text{cm}}\,\sqrt{2\ln 2}.
\end{equation}

\subsubsection*{Marcus--Hush width}

The classical high-temperature Marcus absorption band is a Gaussian in
energy whose variance is \(2\lambda k_B T\), giving
\begin{equation}
\sigma_E = \sqrt{2\,\lambda\,k_B T}.
\end{equation}
With \(k_B T = 0.025693\)\,eV at 298.15\,K, for \(\lambda = 0.30\)\,eV:
\(\sigma_E = 0.12416\)\,eV (FWHM \(= 0.29234\)\,eV).

\medskip
\noindent\textbf{Caveat.} This is the classical (high-temperature) limit.
For organic vibrational baths (\(\sim\)1400\,cm\(^{-1}\)),
\(\hbar\omega \gg k_B T\), so the Marcus--Hush width is a first-order
estimate rather than a quantitative bandshape. For quantitative bandshape
work, supply vibronic (Franck--Condon) spectra through the
experimental/spreadsheet (CSV/XLSX) input branch. The tool is not locked
to Gaussian broadening and processes any tabular absorption spectrum
through the identical descriptor pipeline.

\subsection{Calibration}
\label{sec:calibration}

A calibration block supplied via \texttt{-{}-calibration PATH} (CLI and
library) is applied \emph{before} broadening and does not change the
output path.

The calibration performs three operations, all optional and independent:
\begin{enumerate}
  \item \textbf{Energy map} \(E_{\text{cal}} = aE + b\) (E and b in eV).
        Wavelength is recomputed from the calibrated energy via
        \(\lambda = 1239.841984\,\text{eV\,nm}\,/\,E_{\text{cal}}\) only
        when the map is non-identity (i.e.\ when \(a \ne 1\) or
        \(b \ne 0\)).
  \item \textbf{Oscillator-strength scaling}
        \(f_{\text{cal}} = \alpha f\).
  \item \textbf{Per-band overrides}: optional per-transition
        \texttt{sigma\_ev} or \texttt{reorganization\_ev} that take
        precedence over the global values.
\end{enumerate}

\noindent Guards: \(\alpha > 0\) and every \(E_{\text{cal}} > 0\) are
enforced. The identity calibration (\(a=1\), \(b=0\), \(\alpha=1\), no
band overrides) produces output bit-identical to running without
\texttt{-{}-calibration}.

\paragraph{TOML schema.}
\begin{lstlisting}
[energy]
a = 1.0
b = 0.0

[oscillator]
alpha = 1.0

[[band]]
index = 1            # 1-based excited-state index
sigma_ev = 0.25      # width override (eV)
reorganization_ev = 0.40
\end{lstlisting}

All sections and all keys are optional. Omitted keys take their identity
or default values.

\medskip
\noindent\textbf{Note.} \texttt{-{}-calibration} is available on the CLI
and via the Python library. The HTTP API does not accept calibration files.

% ====================================================================
\section{Use as a Library}
\label{sec:library}

\texttt{overlap\_calculator} is installable as a Python package and
exposes a fully typed public API. Every name below is importable directly
from \texttt{overlap\_calculator}.

\subsection*{Minimal snippet}

\begin{lstlisting}[language=Python]
from overlap_calculator import analyze
from overlap_calculator.models import AnalysisInput

items = [AnalysisInput(
    input_type="theoretical",
    source_path="tddft.out",
    sample_id="dye-1",
)]
results, skipped = analyze(items, default_light_sources=["AM15G"])
\end{lstlisting}

\subsection*{\texttt{analyze} signature}

\begin{lstlisting}[language=Python]
analyze(
    items: list[AnalysisInput],
    *,
    sigma_ev: float = 0.30,
    wl_min: float = 200.0,
    wl_max: float = 800.0,
    num_points: int = 10000,
    concentration_molar: float = 1e-5,
    path_cm: float = 1.0,
    default_light_sources: list[str] | None = None,
    custom_light_source_paths: list[Path] | None = None,
    prefactor_mode: str = "constant",
    sigma_mode: str = "fixed",
    reorganization_ev: float = 0.30,
    temperature_k: float = 298.15,
    calibration: Calibration | None = None,
) -> tuple[list[RunResult], list[AnalysisSkip]]
\end{lstlisting}

\subsection*{Public API names}

\begin{longtable}{@{}>{\RaggedRight}p{5.4cm}>{\RaggedRight}p{10.0cm}@{}}
\toprule
\textbf{Name} & \textbf{Purpose} \\
\midrule
\endhead
\texttt{analyze}                & High-level entry point --- full pipeline \\
\texttt{analyze\_inputs}        & Lower-level batch runner \\
\texttt{export\_results}        & Write tables and plots from a \texttt{results} list \\
\texttt{prepare\_output\_dir}   & Create the output directory tree \\
\idtt{build\_extinction\_spectrum} & Reconstruct \(\varepsilon(\lambda)\) from TD-DFT excited states \\
\texttt{compute\_absorbance}    & Beer--Lambert \(A(\lambda) = \varepsilon c L\) \\
\texttt{compute\_absorptance}   & \(\alpha(\lambda) = 1 - 10^{-A}\) \\
\texttt{shape\_overlap}         & Dimensionless spectral shape comparator \\
\texttt{integrate\_light\_flux} & \(\int I(\lambda)\,d\lambda\) \\
\idtt{integrate\_absorbed\_flux} & \(\int \alpha(\lambda)\,I(\lambda)\,d\lambda\) \\
\texttt{load\_light\_sources}   & Load bundled and custom light sources \\
\texttt{marcus\_hush\_sigma\_ev}& \(\sigma = \sqrt{2\lambda k_B T}\) in eV \\
\texttt{Calibration}            & Pydantic model for the calibration block \\
\texttt{load\_calibration}      & Parse a TOML calibration file \\
\texttt{apply\_calibration}     & Apply a calibration to a list of excited states \\
\idtt{build\_run\_manifest}     & Assemble a provenance manifest dict \\
\idtt{write\_run\_manifest}     & Write \texttt{run\_manifest.json} to disk \\
\texttt{\_\_version\_\_}        & Installed package version string \\
\bottomrule
\end{longtable}

% ====================================================================
\section{Configuration and Environment Variables}
\label{sec:config}

All settings are read via \texttt{pydantic-settings} with the
\texttt{GAUSS\_} prefix; a \texttt{.env} file in the working directory
is honoured.

\begin{longtable}{@{}>{\RaggedRight}p{7.2cm}>{\RaggedRight}p{3.6cm}>{\RaggedRight}p{3.95cm}@{}}
\toprule
\textbf{Variable} & \textbf{Default} & \textbf{Purpose} \\
\midrule
\endhead
\idtt{GAUSS\_LOG\_LEVEL}                & \texttt{INFO}                               & Logging verbosity \\
\idtt{GAUSS\_LOG\_FORMAT}               & \texttt{text}                               & \texttt{text} or \texttt{json} \\
\idtt{GAUSS\_PLOT\_OUTPUTS}             & \texttt{true}                               & Emit TIFF plots \\
\idtt{GAUSS\_PLOT\_DPI}                 & \(400\)                                     & Default TIFF plot resolution \\
\idtt{GAUSS\_RANKING\_OUTPUTS}          & \texttt{true}                               & Emit grouped ranking tables and bar charts \\
\idtt{GAUSS\_SIGMA\_EV}                 & \(0.30\)                                    & Default broadening \(\sigma\) \\
\idtt{GAUSS\_WAVELENGTH\_MIN\_NM}       & \(200.0\)                                   & Grid lower bound, nm \\
\idtt{GAUSS\_WAVELENGTH\_MAX\_NM}       & \(800.0\)                                   & Grid upper bound, nm \\
\idtt{GAUSS\_WAVELENGTH\_POINTS}        & \(10000\)                                   & Number of grid points \\
\idtt{GAUSS\_REFERENCE\_CONCENTRATION\_MOLAR} & \(1\text{e-}5\)                       & Beer--Lambert \(c\) \\
\idtt{GAUSS\_REFERENCE\_PATH\_CM}       & \(1.0\)                                     & Beer--Lambert \(L\) \\
\idtt{GAUSS\_DEFAULT\_LIGHT\_SOURCES}   & \idtt{AM15G,LEDB4,\ldots}                   & Default bundled-source list \\
\idtt{GAUSS\_UPLOAD\_MAX\_MB}           & \(50\)                                      & API upload size cap \\
\idtt{GAUSS\_UPLOAD\_ALLOWED\_EXTS}     & \idtt{.out,.log,.csv,.xlsx,.xls}            & Accepted upload extensions \\
\bottomrule
\end{longtable}

\noindent CLI flags override environment variables when both are
supplied.

% ====================================================================
\section{Troubleshooting}

\begin{description}[style=nextline,leftmargin=1em,labelindent=0pt,font=\normalfont\ttfamily]
  \item[No output generated]
  Inspect \texttt{output/tables/skipped\_inputs.*} first and re-run
  with \texttt{-{}-log-level DEBUG}. A common cause is a \texttt{.out}
  file that is an OPT job, not a TD-DFT job, so it contains no
  \texttt{Excited State} blocks.

  \item[ParseError on an experimental file]
  Confirm one column is a wavelength axis (alias list in \S\ref{sec:input})
  and at least one other column is numeric. Non-numeric /
  date-formatted columns are dropped silently.

  \item[Custom light source is ignored]
  Check that the file has exactly one wavelength column and one
  numeric intensity column. The filename stem becomes the
  \texttt{light\_source\_name}; if that name collides with a bundled
  source the custom file wins.

  \item[CLI options seem outdated]
  After switching branches, re-install the package:
  \texttt{pip install -e .}.

  \item[Unexpected absorbed\_fraction near 0 for a CIE LED source]
  Check that the dye's \(\lambda_\text{max}\) falls within the LED
  emission band. \texttt{shape\_overlap} is often more informative
  than \texttt{absorbed\_fraction} for narrow-band relative sources.

  \item[\texttt{run\_manifest.json} is missing]
  The manifest is written by \texttt{analyze}. It will not be present
  for runs that were aborted before completion. Re-run the pipeline to
  produce it.

  \item[\texttt{git\_commit} is \texttt{null} in \texttt{run\_manifest.json}]
  The package was installed from a wheel or archive rather than from a
  git checkout, so no commit SHA is available. The
  \texttt{software\_version} field still identifies the release.
\end{description}

% ====================================================================
\section{Citation and License}
\label{sec:citation}

If you use \texttt{overlap-calculator} in published work, please
cite the concept DOI below. It resolves to the latest archived
release on Zenodo, so the citation stays valid across versions
without needing an update.

\begin{quote}
Seyitdanl{\i}o{\u{g}}lu, P. \emph{overlap-calculator: Batch
spectral-overlap analysis between TD-DFT / experimental absorption and
reference light sources.} Zenodo.
\url{https://doi.org/10.5281/zenodo.19944515}
\end{quote}

BibTeX:
\begin{lstlisting}
@software{seyitdanlioglu_overlap_calculator,
  author    = {Seyitdanl{\i}o{\u{g}}lu, P{\i}nar},
  title     = {{overlap-calculator}: Batch spectral-overlap analysis
               between TD-DFT / experimental absorption and reference
               light sources},
  year      = {2026},
  publisher = {Zenodo},
  doi       = {10.5281/zenodo.19944515},
  url       = {https://doi.org/10.5281/zenodo.19944515}
}
\end{lstlisting}

\subsection*{Acknowledgement of third-party data}

The bundled CIE LED illuminants (\texttt{LEDB1}, \texttt{LEDB2},
\texttt{LEDB3}, \texttt{LEDB4}) and fluorescent illuminant
(\texttt{CIEFL10}) are official \textbf{CIE} data sets derived from
\emph{CIE 015:2018 Colorimetry, 4th Edition}. The \texttt{AM15G}
spectrum is the ASTM G-173-03 AM1.5 global reference distributed by
the U.S.\ National Renewable Energy Laboratory (NREL). The bundled
experimental absorbance workbook
(\texttt{input/files/organic\_uvvis\_photochemcad\_dataset.xlsx}) is
derived from the \textbf{PhotochemCAD} spectral database. If you use
\texttt{overlap-calculator} in published work, please cite the data
sets alongside the software:

\begin{quote}
CIE (2018). \emph{Relative spectral power distributions of
illuminants representing typical LED lamps, 1\,nm spacing.}
International Commission on Illumination, Vienna, AT.
DOI: \href{https://doi.org/10.25039/CIE.DS.dhcw57sd}{10.25039/CIE.DS.dhcw57sd}
\end{quote}

\begin{quote}
CIE (2018). \emph{Relative spectral power distributions of
illuminants representing typical fluorescent lamps, 1\,nm wavelength
steps.} International Commission on Illumination, Vienna, AT.
DOI: \href{https://doi.org/10.25039/CIE.DS.54hy6srn}{10.25039/CIE.DS.54hy6srn}
\end{quote}

\begin{quote}
CIE (2018). \emph{CIE 015:2018 Colorimetry, 4th Edition.}
International Commission on Illumination, Vienna, AT.\\
\url{https://cie.co.at/publications/colorimetry-4th-edition/}
\end{quote}

\begin{quote}
NREL. \emph{Reference Air Mass 1.5 Spectra (ASTM G-173-03).} U.S.\
National Renewable Energy Laboratory.
\url{https://www.nrel.gov/grid/solar-resource/spectra-am1.5.html}
\end{quote}

\begin{quote}
Taniguchi, M., \& Lindsey, J.~S. (2018). Database of absorption and
fluorescence spectra of {>}300 common compounds for use in
PhotochemCAD. \emph{Photochemistry and Photobiology}, \textbf{94}(2),
290--327.
DOI: \href{https://doi.org/10.1111/php.12860}{10.1111/php.12860}
\end{quote}

BibTeX:
\begin{lstlisting}
@techreport{cie_led_illuminants_2018,
  author      = {{International Commission on Illumination}},
  title       = {Relative spectral power distributions of illuminants
                 representing typical {LED} lamps, 1\,nm spacing},
  institution = {CIE Central Bureau},
  address     = {Vienna, AT},
  year        = {2018},
  doi         = {10.25039/CIE.DS.dhcw57sd},
  url         = {https://doi.org/10.25039/CIE.DS.dhcw57sd}
}

@techreport{cie_fl_illuminants_2018,
  author      = {{International Commission on Illumination}},
  title       = {Relative spectral power distributions of illuminants
                 representing typical fluorescent lamps, 1\,nm wavelength steps},
  institution = {CIE Central Bureau},
  address     = {Vienna, AT},
  year        = {2018},
  doi         = {10.25039/CIE.DS.54hy6srn},
  url         = {https://doi.org/10.25039/CIE.DS.54hy6srn}
}

@techreport{cie_015_2018,
  author      = {{International Commission on Illumination}},
  title       = {{CIE} 015:2018 Colorimetry, 4th Edition},
  institution = {CIE Central Bureau},
  address     = {Vienna, AT},
  year        = {2018},
  isbn        = {978-3-902842-13-8}
}

@misc{nrel_am15g,
  author       = {{National Renewable Energy Laboratory}},
  title        = {Reference Air Mass 1.5 Spectra ({ASTM} {G-173-03})},
  howpublished = {\url{https://www.nrel.gov/grid/solar-resource/spectra-am1.5.html}},
  institution  = {U.S. Department of Energy}
}

@article{taniguchi_lindsey_photochemcad_2018,
  author  = {Taniguchi, Masahiko and Lindsey, Jonathan S.},
  title   = {Database of Absorption and Fluorescence Spectra of
             {>}300 Common Compounds for Use in {PhotochemCAD}},
  journal = {Photochemistry and Photobiology},
  volume  = {94},
  number  = {2},
  pages   = {290--327},
  year    = {2018},
  doi     = {10.1111/php.12860},
  url     = {https://doi.org/10.1111/php.12860}
}
\end{lstlisting}

Please consult the CIE at \url{https://cie.co.at/} for the
authoritative reference tables and the latest errata.

\texttt{overlap-calculator} is released under the \textbf{GNU General
Public License v3.0 or later} (see \texttt{LICENSE} in the source
distribution).

\end{document}
